% Style based on the supplied Waseem_Generic_AI_ML_GenAI_AgenticAI.tex.
% Python replaces template tokens. Edit layout here, content in JSON/portfolio JS.
\documentclass[9pt,letterpaper]{extarticle}
\usepackage[margin=0.55in]{geometry}
\usepackage[T1]{fontenc}
\usepackage[utf8]{inputenc}
\usepackage{lmodern}
\usepackage{microtype}
\usepackage{xcolor}
\usepackage{enumitem}
\usepackage{titlesec}
\usepackage[hidelinks]{hyperref}
\usepackage{tabularx,array,parskip,ragged2e,needspace}
\pagestyle{empty}
\setlength{\parindent}{0pt}
\setlength{\parskip}{1pt}
\definecolor{accent}{HTML}{155B8A}
\titleformat{\section}{\large\bfseries\color{accent}}{}{0pt}{}[\vspace{0.5pt}\color{accent}\titlerule]
\titlespacing*{\section}{0pt}{3pt}{2pt}
\setlist[itemize]{leftmargin=1.15em,itemsep=0.8pt,topsep=1pt,parsep=0pt,partopsep=0pt}
\newcommand{\role}[4]{%
  \Needspace{4\baselineskip}\textbf{#1}\hfill\textbf{#2}\\[-1pt]
  \textit{#3}\hfill\textit{#4}\par}
\newcommand{\entry}[2]{\Needspace{4\baselineskip}\textbf{#1}\hfill\textit{#2}\par}
\newcommand{\skillrow}[2]{\textbf{#1:} & #2\\}
\hypersetup{pdftitle={Dr Waseem Raza, PhD - Healthcare AI \& Wearable Intelligence},pdfauthor={Dr Waseem Raza, PhD}}
\begin{document}
\begin{center}
{\LARGE\bfseries Dr Waseem Raza, PhD}\\[2pt]
{\large\bfseries Machine Learning Engineer \textbar{} Healthcare AI \& Wearable Intelligence}\\[1pt]
{\small Healthcare AI \textbar{} Wearable Sensing \textbar{} Deep Learning \textbar{} Reliable ML}\\[3pt]
{\small New Jersey, United States \quad 732-803-9703 \quad \href{mailto:waseemraza844@gmail.com}{waseemraza844@gmail.com} \quad \href{https://www.linkedin.com/in/waseem-raza-phd-9a737858/}{LinkedIn} \quad \href{https://github.com/WaseemRaza844}{GitHub} \quad \href{https://scholar.google.com/citations?user=d58Drr0AAAAJ\&hl=en}{Scholar} \quad \href{https://waseemraza844.github.io/portfolio/}{Portfolio}}
\end{center}
\vspace{-5pt}
\Needspace{3\baselineskip}
\section{Professional Summary}
Machine Learning Engineer and researcher with a PhD in Electrical and Computer Engineering and experience applying deep learning, time-series modeling, and reliable AI to wearable sensing and real-world health-monitoring problems. Research includes continuous medication-taking recognition from wrist-worn inertial sensors, personalization, domain adaptation, model evaluation, and cross-user robustness, complemented by broader production ML and cloud-data engineering experience.
\Needspace{3\baselineskip}
\section{Technical Expertise, Leadership, and Professional Strengths}
\Needspace{2\baselineskip}\noindent{\small\textbf{Healthcare AI \& Wearables:} Wearable sensing, medication-adherence monitoring, activity recognition, smartwatch IMU data, human-centered sensing, and health-data applications}\par
\Needspace{2\baselineskip}\noindent{\small\textbf{Deep Learning:} CNNs, LSTMs, attention, transfer learning, domain adaptation, representation learning, classification, and sequence modeling}\par
\Needspace{2\baselineskip}\noindent{\small\textbf{Time-Series \& Sensor ML:} Windowing, segmentation, feature engineering, temporal modeling, cross-user evaluation, personalization, and continuous-event detection}\par
\Needspace{2\baselineskip}\noindent{\small\textbf{Reliable \& Personalized ML:} Distribution shift, few-shot adaptation, fairness-aware evaluation, robustness analysis, LOSO validation, and error analysis}\par
\Needspace{2\baselineskip}\noindent{\small\textbf{Data \& Engineering:} Python, PyTorch, TensorFlow, scikit-learn, Pandas, NumPy, cloud analytics, data pipelines, APIs, and reproducible experimentation}\par
\Needspace{2\baselineskip}\noindent{\small\textbf{Research \& Evaluation:} Experimental design, statistical analysis, sensitivity/specificity, precision/recall/F1, cross-validation, visualization, and technical communication}\par
\Needspace{4\baselineskip}\textbf{Leadership and Professional Strengths}\par
\begin{itemize}
\item \textbf{Ownership and accountability:} Lead AI/ML initiatives through requirements analysis, data preparation, model development, validation, stakeholder review, and deployment planning.
\item \textbf{Technical leadership:} Define technical approaches, evaluate alternatives, coordinate experimentation, and communicate findings to guide projects from ambiguous requirements to measurable outcomes.
\item \textbf{Cross-functional collaboration:} Work with software engineers, network specialists, researchers, geologists, and public-safety stakeholders to translate domain requirements into practical data and AI solutions.
\item \textbf{Continuous learning:} Apply ongoing study in generative AI, agentic AI, RAG, LLM evaluation, and cloud-native platforms through hands-on projects and professional coursework.
\item \textbf{Reliability and commitment:} Maintain reproducible workflows, careful validation, and attention to detail throughout experimentation and delivery.
\item \textbf{Adaptability and problem solving:} Combine critical thinking, systematic experimentation, and data-driven analysis to address unfamiliar technologies and complex operational problems.
\end{itemize}
\Needspace{3\baselineskip}
\section{Professional Experience}
\role{Machine Learning Engineer}{Feb. 2025-Present}{AT\&T, contractor via Innovatix Inc.}{Bedminster, NJ}
\begin{itemize}\interlinepenalty=10000
\item Own end-to-end development, testing, and operational integration of an AI-driven diagnostic system for critical E911 service-assurance workflows, improving automated root-cause classification accuracy from 20\% to 80\%.
\item Designed domain-adaptive neural models and reusable data-processing workflows that generalize across heterogeneous operational domains and reduce performance degradation caused by distribution shifts.
\item Developed multimodal AI prototypes combining work-order text, location information, network context, and operational signals to support GenAI-assisted diagnostics for VoNR and 5G service assurance.
\item Created model-validation, error-analysis, and failure-mode assessment procedures to compare performance across domains and support reliable deployment in high-impact operational environments.
\item Partnered with network engineers, software teams, and public-safety stakeholders to translate operational requirements into validated AI workflows and production-oriented data assets.
\end{itemize}
\role{Data Scientist and AI/ML Engineer}{Sep. 2024-Dec. 2024}{Impac Exploration Services}{Weatherford, OK}
\begin{itemize}\interlinepenalty=10000
\item Built an LLM-enabled analytical pipeline for extracting and interpreting information from geochemical well data.
\item Developed scalable ingestion, cleansing, transformation, and validation workflows using Python, Azure Databricks, Azure Functions, and Synapse Analytics.
\item Automated ETL processing for large exploration datasets, reducing data-preparation time by 40\% and improving reproducibility.
\item Implemented cross-validation, data-quality, and consistency checks that improved reliability of analytical outputs by 20\%.
\item Collaborated with geologists to convert domain requirements and technical information into AI-ready datasets and actionable analytical outputs.
\end{itemize}
\role{Graduate Research Assistant}{Sep. 2019-Sep. 2024}{AI4Networks, University of Oklahoma}{Tulsa, OK}
\begin{itemize}\interlinepenalty=10000
\item Designed reusable Python simulation and data-generation platforms for large-scale experimentation, including a 3GPP-compliant network simulator and an AI-enabled 5G testbed.
\item Built transfer-learning, adversarial-learning, GAN, autoencoder, and reinforcement-learning pipelines for sparse-data and cross-domain modeling.
\item Developed robustness evaluation frameworks for models operating under distribution shifts, missing data, and limited labeled samples; led 100+ experiments from problem formulation through validation.
\item Developed scalable self-healing and outage-management frameworks that improved modeled network resilience by 30\%.
\item Published peer-reviewed research, presented technical results internationally, and mentored junior researchers in ML implementation, experimental design, and technical writing.
\end{itemize}
\Needspace{3\baselineskip}
\section{Education}
\entry{PhD, Electrical and Computer Engineering}{09/2019-07/2025}
University of Oklahoma, Norman, Oklahoma, USA\par
\entry{MSc, Telecommunication Engineering}{08/2014-08/2016}
University of Engineering and Technology, Taxila, Pakistan\par
\entry{BSc, Telecommunication Engineering}{11/2010-06/2014}
University of Engineering and Technology, Taxila, Pakistan\par
\Needspace{3\baselineskip}
\section{Selected Projects}
\entry{Wearable Intelligence for Health Monitoring}{2024-2026}
\textit{Research Project}\par
\begin{itemize}\interlinepenalty=10000
\item Deep-learning methods for continuous recognition of medication and dietary activities from wrist-worn inertial sensors.
\end{itemize}
{\small\textit{Tools \& methods:} Wearable AI, Time Series, Deep Learning, Smartwatch Sensing\par}
\entry{Cloud Computer Vision}{2025}
\textit{Guided Project}\par
\begin{itemize}\interlinepenalty=10000
\item Real-time object detection, vehicle counting, and image-classification workflows using cloud-native AI services.
\end{itemize}
{\small\textit{Tools \& methods:} Computer Vision, Streaming, Google Cloud AutoML, Python\par}
\Needspace{3\baselineskip}
\section{Certifications and Continuing Education}
\Needspace{3\baselineskip}\noindent\textbf{1. Deep Learning for Healthcare Specialization} {\small --- \textcolor{blue}{\underline{\href{https://www.coursera.org/specializations/deep-learning-healthcare}{University of Illinois Urbana-Champaign / Coursera}}} \textbar{} Completed (0/3) courses}\par
\textbf{Focus:} An advanced specialization covering health-data processing, neural-network methods, and the development of deep-learning solutions for real-world medical and healthcare applications.\par
{\small\textbf{Tools \& Skills:} Deep Learning, Healthcare AI, Health Informatics, Health Data Processing, Artificial Neural Networks, Convolutional Neural Networks, Recurrent Neural Networks, Autoencoders, Generative Models, Medical Image Analysis, Supervised Learning, Unsupervised Learning, Predictive Modeling, Model Evaluation, PyTorch, Jupyter Notebooks\par}
\vspace{2pt}
\Needspace{3\baselineskip}\noindent\textbf{2. Machine Learning} {\small --- Stanford University / Coursera \textbar{} Completed}\par
\textbf{Focus:} Core supervised and unsupervised learning, model evaluation, optimization, and practical ML system development.\par
{\small\textbf{Tools \& Skills:} Supervised Learning, Unsupervised Learning, Model Evaluation, Optimization\par}
\vspace{2pt}
\Needspace{3\baselineskip}\noindent\textbf{3. Google Cloud Data Analytics Professional Certificate} {\small --- Google Cloud / Coursera \textbar{} Completed}\par
\textbf{Focus:} Cloud data transformation, analysis, visualization, and decision support using Google Cloud tools.\par
{\small\textbf{Tools \& Skills:} SQL, BigQuery, Data Analytics, Visualization\par}
\vspace{2pt}
\Needspace{3\baselineskip}\noindent\textbf{4. IBM Generative AI Engineering Professional Certificate} {\small --- IBM / Coursera \textbar{} Completed (13/16) courses}\par
\textbf{Focus:} Sixteen-course professional certificate developing job-ready skills across AI foundations, Python application development, machine learning, deep learning, NLP, transformer architectures, LLM fine-tuning, RAG, LangChain, and AI agents.\par
{\small\textbf{Tools \& Skills:} Generative AI Engineering, Large Language Models, Transformer Architectures, Prompt Engineering, RAG, LangChain, AI Agents, Hugging Face, PyTorch, Keras, Scikit-learn, Python, Flask, Gradio, PEFT, LoRA, QLoRA, RLHF, DPO, PPO, NLP, Machine Learning, Deep Learning, Vector Databases\par}
\vspace{2pt}
\Needspace{3\baselineskip}
\section{Selected Publications}
\entry{Real-World Medication Adherence Monitoring Using Hybrid Neural Networks and Smartwatches}{2026}
\textit{HealthINF  \textbar{}  Accepted}\par
\end{document}
